\documentclass[11pt,a4paper]{article}

\usepackage[margin=1in]{geometry}
\usepackage{booktabs}
\usepackage{hyperref}
\usepackage{xcolor}
\usepackage{enumitem}

\definecolor{revisionblue}{RGB}{0,102,204}
\hypersetup{
    colorlinks=true,
    linkcolor=revisionblue,
    urlcolor=revisionblue
}

\begin{document}

\begin{center}
    {\Large\bfseries Response to Reviewers}\\[0.5em]
    {\large IEEE Transactions on Multimedia}\\[1em]
    {\bfseries Manuscript ID: \#8159}\\[0.3em]
    {\bfseries Title: Texture Resonance Retrieval for Retrieval-Grounded Editable Audio Effect Preset Selection}\\[1.2em]
\end{center}

\noindent Dear Editor and Reviewers,

\medskip

We thank the reviewers for the detailed critiques. The central lesson from the reviews was that our earlier package mixed a verified retrieval benchmark with broader system claims that were not supported to the same standard. In the revision, we therefore prioritized \emph{evidence hardening} over scope expansion. The revised manuscript now makes fewer claims, but each retained claim is tied to the verified P0 experiment package or a canonical repository artifact.

\section*{Canonical Evidence Sources}

The revised submission uses the following repository artifacts as the sole experimental sources of truth:

{\footnotesize
\begin{itemize}[leftmargin=2em]
    \item \path{Experiments/tmm/protocolA_split_audit.md}
    \item \path{Experiments/tmm/protocolA_audio_grouped_leakage_report.json}
    \item \path{Experiments/AblationStudies/protocolA_audio_grouped_objective_stats.md}
    \item \path{Experiments/AblationStudies/protocolA_audio_grouped_per_query_metrics.csv}
    \item \path{Experiments/AblationStudies/outputs/p0_e2/e2_summary.csv}
    \item \path{Experiments/AblationStudies/outputs/p0_e2/e2_stats.json}
    \item \path{Experiments/AblationStudies/outputs/p0_e3_shared_e2_test/e3_overall_summary.csv}
    \item \path{Experiments/AblationStudies/outputs/p0_e3_shared_e2_test/e3_summary.csv}
    \item \path{Experiments/E5_Ablations/outputs/p0_near_dup/p0_near_dup_sensitivity.csv}
    \item \path{Experiments/E2_SOTABaselines/outputs/p0_mlp_regressor/mlp_regressor_results.json}
    \item \path{Experiments/mushura/results_report.md}
\end{itemize}
}

We do not use the legacy retrieval report, the deprecated Protocol-B outputs, or the repository's exploratory personalization/memory code as formal evidence in this revision.

\section*{Summary of Revision Actions}

\begin{enumerate}[leftmargin=2em]
    \item \textbf{Scope narrowed to retrieval-grounded editable control.} The manuscript now centers on Texture Resonance Retrieval (TRR) as a retrieval representation for executable preset control. Personalization, AEM memory, and related extensions are explicitly presented as exploratory rather than validated core contributions.

    \item \textbf{Objective claims upgraded to the P0 Protocol-A baseline matrix.} A post-hoc audit of the legacy song-name-based Protocol-A list found 16 cross-split shared resolved audio paths. We retain this audit for split transparency and now use the P0 204-query Protocol-A package as the main baseline matrix. The revised comparison includes TRR, Wav2Vec, FeatureNN, CLAP, PaSST, and PANNs on the shared 204-query / 1063-KB setup.

    \item \textbf{Per-query diagnostics added.} We generated a derived analysis directly from the canonical Protocol-A per-query CSV for the stricter split, including win/loss/tie summaries, signed-delta diagnostics, and representative success/failure cases. These diagnostics are now cited in the supplementary material so that the main result is not supported by mean values alone.

    \item \textbf{TRR mechanism ablations added.} We added a controlled same-backbone ablation on the P0 split comparing Wav2Vec2 mean pooling, full unprojected Gram aggregation, single-layer projected Gram, projection dimensions, and with/without L2 normalization. The result supports a bounded mechanism claim: second-order Gram aggregation improves over same-backbone mean pooling and L2 normalization is important, while layer and projection choices remain meaningful hyperparameters.

    \item \textbf{Stronger baselines and boundary baselines added.} We now include CLAP, PaSST, and PANNs in the P0 Protocol-A evidence chain. TRR remains strongest on normalized L2, Acc@0.1, Recall, Cosine, and Module among retrieval methods. We also add a diagnostic direct-regression boundary baseline using mean-pooled Wav2Vec2 features and a two-layer MLP trained only on the P0 knowledge base. Because that regressor predicts dense continuous parameter leaves rather than retrieved executable presets, we report it as a boundary condition rather than as a replacement for the retrieval leaderboard. Learned reranking remains future work.

    \item \textbf{Protocol-C repositioned as diagnostic robustness evidence.} The older Protocol-C diagnostic used synthetic embedding-space degradation and remains only a boundary-condition artifact. The P0 Protocol-C result applies AWGN, MP3, reverb, and truncation to raw audio before TRR re-encoding, reuses the same 204-query split as E2, and compares fixed fusion against adaptive fusion. We additionally report current-code audio-only, pure-text, and parameter-oracle upper-bound rows as an execution-path audit. The revised manuscript no longer treats adaptive fusion as proof of a standalone main contribution or superiority over a pure-text boundary.

    \item \textbf{Listening-study terminology and statistics corrected.} The study is consistently described as a multiple-stimulus listening test with a hidden reference and no explicit low-quality anchor. The revision reports the canonical nonparametric statistics already present in the repository: Friedman omnibus tests, Holm-corrected Wilcoxon signed-rank post-hoc tests, and rank-biserial effect sizes, together with an explicit limitation that the logs do not preserve device metadata, participant expertise, or the provenance/time budget of the manual baseline.

    \item \textbf{Submission package synchronized.} The manuscript, supplement, response documents, and checklist were edited to remove stale claims about PaSST/PANNs, real-audio Protocol-C, near-duplicate placeholders, normalized metrics, and AEM learning curves.
\end{enumerate}

\section*{Responses to the Main Reviewer Themes}

\subsection*{1. Novelty and scope}

We agree that the strongest supported contribution is not a full personalized agent system, but a retrieval-grounded parameter-control formulation centered on TRR. The revised manuscript therefore frames the contribution around texture-aware retrieval for editable preset control and avoids claims that would require additional evidence beyond the current benchmark.

\subsection*{2. Benchmark scale, split transparency, and statistical reporting}

The objective results are now anchored to the shared P0 Protocol-A benchmark with $N_{\mathrm{test}}=204$ and $N_{\mathrm{kb}}=1063$. The revised manuscript and supplementary material report:

\begin{itemize}[leftmargin=2em]
    \item the legacy split audit (16 shared resolved audio paths across split);
    \item the stricter anti-leakage split as audit context;
    \item the P0 baseline matrix including Wav2Vec, FeatureNN, CLAP, PaSST, and PANNs;
    \item method means including L2, Norm.L2, Acc@0.1, Recall, Cosine, and Module;
    \item Wilcoxon tests on per-query Norm.L2 for TRR versus each retrieval baseline;
    \item Holm-corrected decisions and effect-size caveats.
\end{itemize}

\subsection*{3. Stronger baselines}

We now include \emph{CLAP, PaSST, and PANNs} as stronger retrieval baselines in the P0 Protocol-A package. The revised manuscript therefore makes the narrower and better-supported statement that TRR is the strongest among the currently evaluated retrieval baselines on the shared P0 benchmark, including Wav2Vec, FeatureNN, CLAP, PaSST, and PANNs. We also include a diagnostic direct-regression MLP boundary baseline, but we do not claim to have completed an executable direct-parameter controller or learned reranking baseline.

\subsection*{3a. TRR mechanism evidence}

We agree that the previous version did not sufficiently isolate why TRR works. We therefore added a controlled same-backbone mechanism ablation. The new supplementary table compares Wav2Vec2 mean pooling, full unprojected Gram aggregation, single-layer Gram variants for layers 4/5/6, projected multi-layer Gram variants with dimensions 16/32/64/128, and the same multi-layer setting without L2 normalization. The result supports a narrower mechanism claim: second-order Gram aggregation improves over same-backbone mean pooling, and L2 normalization is important. It also shows that the exact layer and projection choice matters, so the revised paper frames TRR as an effective texture-aware retrieval prior rather than a fully optimized universal representation.

\subsection*{4. Metric limitations}

We also do \emph{not} claim to have replaced the metric system with a normalized or categorical-aware suite. Instead, the revision clarifies the exact operational metric definitions and their limitations:

\begin{itemize}[leftmargin=2em]
    \item flattened numeric leaves only;
    \item missing numeric keys treated as zero;
    \item no per-parameter normalization;
    \item within-protocol interpretation only.
\end{itemize}

These limitations are now stated explicitly in both the main paper and the supplementary material.

\subsection*{5. Protocol-C realism}

We agree with the reviewers that synthetic embedding-space degradation is not equivalent to real audio corruption. The earlier diagnostic Protocol-C remains limited for that reason. The P0 Protocol-C package addresses this gap by applying real degradations to raw audio before re-encoding:

\begin{itemize}[leftmargin=2em]
    \item AWGN at 20/10/5 dB;
    \item MP3 at 128/64/32 kbps;
    \item reverb with RT60 0.6/1.0;
    \item truncation to 50\%/30\%.
\end{itemize}

The retained claim is not that adaptive fusion dominates every condition or improves over a pure-text boundary. We added a boundary audit with audio-only, pure-text, and parameter-oracle upper-bound rows and revised the manuscript to treat Protocol-C as a diagnostic robustness analysis. The supported claim is narrower: audio-heavy fusion is fragile under realistic degradation, and retrieval stability depends strongly on modality weighting.

\subsection*{6. Listening-study rigor}

The revised package now surfaces the canonical listening-study statistics already present in the repository:

\begin{itemize}[leftmargin=2em]
    \item 26 participants and 910 ratings;
    \item Friedman omnibus tests for Trials 2--5 and Trials 6--10;
    \item Holm-corrected Wilcoxon signed-rank post-hoc tests;
    \item rank-biserial correlation ($rbc$);
    \item explicit statement that the raw \texttt{HCAP} label corresponds to the TRR-based system described in the paper.
\end{itemize}

We also state clearly that the TRR-based system and MusicGen are statistically indistinguishable in Trials 6--10 under the reported post-hoc test, so this block is interpreted as parity rather than superiority.

\subsection*{7. Scope mismatch between paper and repository}

This concern was valid in earlier drafts. The revision addresses it by making the repository's exploratory modules non-central:

\begin{itemize}[leftmargin=2em]
    \item AEM / personalization is described as exploratory only;
    \item deprecated Protocol-B results are removed from the primary evidence path;
    \item legacy reports are no longer used as formal evidence;
    \item the main architecture figure now depicts only the verified retrieval-grounded pipeline.
\end{itemize}

\section*{Closing Statement}

We appreciate the reviewers' insistence on a tighter evidence chain. The revised submission is intentionally more conservative than earlier drafts: it makes fewer claims, but each retained claim is now tied to a canonical artifact with clearer statistical support and clearer scope boundaries.

\medskip

\noindent Sincerely,\\
The Authors

\end{document}
